\documentclass[10pt]{letter}
\usepackage[margin=0.82in]{geometry}
\usepackage{microtype}
\signature{Z.-B. Zhou\\Department of Physics, INSTITUTION\\EMAIL}
\address{Department of Physics\\INSTITUTION\\POSTAL ADDRESS}
\begin{document}
\begin{letter}{Editors\\Physical Review Letters}
\opening{Dear Editors,}

Please consider our manuscript, ``Winding-Free Exact Diagonalization of
Quantum Spin Ice,'' for publication in \emph{Physical Review Letters}.

The recent two-peak heat capacity of Ce$_2$Hf$_2$O$_7$ creates a concrete
need to connect high-temperature exchange fits to the low-energy coherent
ice dynamics. In the relevant sign-problem regime, exact diagonalization is
one of the few unbiased dynamical tools, but its standard 32-site FCC torus
supports a four-spin exchange that closes only through a periodic boundary.
Because this artifact enters at second order while the physical
pyrochlore-hexagon ring exchange enters at third order, it dominates
precisely where the quantum-spin-ice expansion is best controlled.

We derive a gauge-invariant transported-dipole label for complete virtual
processes and show that its zero-character projection becomes a simple
polarization-parity mask on the downfolded ice Hamiltonian. On the complete
2,970-state FCC-32 ice manifold, explicit path filtering and the mask agree
entrywise through third order. The protocol removes all wrapping four- and
six-site processes, retains every contractible hexagon with unchanged
amplitude, restores cubic scaling with the contractible-ring energy, and
reorganizes the four-sublattice-traced longitudinal pseudospin structure
factor at all eight allowed cluster momenta.

For Ce$_2$Hf$_2$O$_7$, the published seventh-order NLC parameter set places
the clean FCC-32 ring peak near 7.1 mK, well below the observed 25 mK
feature. We derive a high-temperature-moment-constrained parameter candidate
that places the ring peak at 25 mK while preserving the leading exchange
second moment and transverse anisotropy of the published set. We clearly
identify this candidate as a target for a new NLC evaluation rather than as
a completed linked-cluster refit.

We believe the manuscript is appropriate for PRL because it provides an
operator-level correction to a parametrically large finite-size error in a
sign-problem setting, validates the correction rigorously on FCC-32, and
turns a newly observed material scale into a concrete constraint on future
joint linked-cluster and winding-free calculations.

This manuscript is original, is not under consideration elsewhere, and has
been approved by all authors. [Add related-manuscript and preprint statement,
if applicable.]

\closing{Sincerely,}
\end{letter}
\end{document}
